\subsection{A rank-one curve with no nonzero three-primary theta point}
\label{sec:no-three-primary-theta}

\paragraph{Finding A20 (all powers of three fail on the degree-thirteen example).}
For the genus-two rank-one curve of A09,
\[
 C/\mathbf F_2:\quad Y^2+X^2Y=X^5+X,
\]
its Raynaud theta divisor satisfies
\[
       \Theta_C(\overline{\mathbf F}_2)
          \cap J_C[3^\infty](\overline{\mathbf F}_2)=\{0\}.
 \tag{TP.1}
\]
Thus allowing arbitrary powers of the single prime $3$ does not
repair the universal positive-rank incidence assertion. This does
not obstruct point separation on this curve: A09 gives its exact
order-$13$ theta point and separating cover.

\begin{proof}
The curve and its Frobenius twist have the same equation over
$\mathbf F_2$. Finding A09 gives its Frobenius characteristic
polynomial on every odd-adic Tate module:
\[
                         P(T)=T^4+T^3+2T+4.
 \tag{TP.2}
\]
Let $\overline P$ be its reduction modulo $3$. Exact polynomial
division in $\mathbf F_3[T]$ gives
\[
                       T^{20}-1=\overline P(T)Q(T),
 \tag{TP.3}
\]
where
\[
\begin{aligned}
 Q(T)={}&T^{16}-T^{15}+T^{14}+T^{12}+T^{11}+T^{10}\\
        &-T^6+T^5-T^4-T^2-T-1.
\end{aligned}
\]
Cayley--Hamilton applied to the action of the $2$-power Frobenius
$\Phi$ on $J_C[3]$ therefore proves $\Phi^{20}=1$ there.
All $3^4$ geometric points of $J_C[3]$ are consequently rational
over $\mathbf F_{2^{20}}$.

The determinant formula for the rational points of a Jacobian
gives
\[
\begin{aligned}
 \#J_C(\mathbf F_{2^{20}})
   &=\det(1-\Phi^{20})
     =\operatorname{Res}(P(T),T^{20}-1)\\
   &=1099812538656
     =3^4\cdot13577932576,
\end{aligned}
 \tag{TP.4}
\]
and $13577932576\equiv1\pmod3$.
The resultant has the displayed sign because $P$ has even degree.
One exact computation uses the integral companion matrix
\[
 C_P=\begin{pmatrix}
 0&0&0&-4\\
 1&0&0&-2\\
 0&1&0&0\\
 0&0&1&-1
 \end{pmatrix}
\]
and evaluates $\det(I-C_P^{20})$; an independent polynomial
resultant calculation gives the same integer.
The determinant formula and Frobenius convention are those of
\cite[Theorem~19.1(b), pp.~40--41]{MilneAV}.

The rational $3$-primary subgroup in (TP.4) has exactly $3^4$
points and already contains all of $J_C[3]$, so it equals
$J_C[3]$ and is killed by $3$.
Apply A19 with the prime set $S=\{3\}$ and the Frobenius exponent
$d=20$ established in (TP.3). It puts every geometric
$3$-primary theta point in $J_C(\mathbf F_{2^{20}})$.
Every nonzero such point would thus have exact order $3$.
Finding A10 excludes order-three theta points on every nonordinary
genus-two curve, and A09 gives $f_C=1$. This is a contradiction.
Finally the origin is on theta, since in A09 the locally exact
differential line bundle is $\mathcal O_C(P)$ for the point at
infinity. This proves equality in (TP.1).
\end{proof}

The standard-library certificate \texttt{check\_three\_primary.py}
performs the companion-matrix and modular polynomial computations.
The independently written SymPy certificate
\texttt{verify\_three\_primary.py} uses a polynomial resultant
and imports no producer code. These verify the finite arithmetic
in (TP.3)--(TP.4); the geometric exclusion also uses the separately
proved A19 and A10.
